\documentclass[10pt,a4paper]{amsart}
\usepackage[margin=19mm]{geometry}
\usepackage{amsmath,amssymb,mathtools}
\usepackage[scr=rsfs,cal=euler]{mathalfa}
\usepackage{xfrac}
\usepackage[unicode,hidelinks,bookmarks=false]{hyperref}
\newcommand{\ie}{i.e.\ }
\newcommand{\cf}{cf.\ }
\newcommand{\resp}{respectively\ }
\newcommand{\Subsection}{\subsection}
\providecommand{\nobreakdash}{\nobreak\hbox{-}}
\setlength{\emergencystretch}{2em}
\multlinegap0pt
\newcommand{\skpt}{\par\smallskip}
\providecommand{\eg}{e.g.\ }
\usepackage{enumerate}
\usepackage{pgfplots}
\usepackage{mathrsfs}
\usepackage{xcolor}
\usepackage{tikz-cd}
\usepackage{scalerel}
\usepackage{circuitikz}

%%%%%%%%%%%%%%%%%%%%%%%%%%%%%%%%%%%%%%%%%
\definecolor{lblue}{rgb}{0.0, 0.55, 1.0}
\definecolor{dred}{rgb}{1.0, 0.0, 0.3}
\definecolor{byzant}{rgb}{0.74, 0.2, 0.64}
\definecolor{pomme}{rgb}{0.16, 0.67, 0.53}
\definecolor{aqua}{rgb}{0.0, 1.0, 1.0}
\definecolor{azure}{rgb}{0.0, 0.5, 1.0}
\definecolor{carnelian}{rgb}{0.7, 0.11, 0.11}
\definecolor{ceruleanblue}{rgb}{0.16, 0.32, 0.75}
\definecolor{pgreen}{rgb}{0.0, 0.65, 0.58}
\definecolor{pblue}{rgb}{0.11, 0.22, 0.73}
\definecolor{pear}{rgb}{0.82, 0.89, 0.19}
\definecolor{pyellow}{rgb}{1.0, 0.97, 0.12}
\definecolor{lgreen}{rgb}{0.75, 1.0, 0.0}
\definecolor{brown}{rgb}{0.6, 0.4, 0.08}
\definecolor{ggray}{rgb}{0.27, 0.35, 0.27}

\definecolor{lblue}{rgb}{0.0, 0.55, 1.0}
\definecolor{dred}{rgb}{1.0, 0.0, 0.3}
\definecolor{byzant}{rgb}{0.74, 0.2, 0.64}
\definecolor{pomme}{rgb}{0.16, 0.67, 0.53}
\definecolor{aqua}{rgb}{0.0, 1.0, 1.0}
\definecolor{azure}{rgb}{0.0, 0.5, 1.0}
\definecolor{carnelian}{rgb}{0.7, 0.11, 0.11}
\definecolor{ceruleanblue}{rgb}{0.16, 0.32, 0.75}
\definecolor{pgreen}{rgb}{0.0, 0.65, 0.58}
\definecolor{pblue}{rgb}{0.11, 0.22, 0.73}
\definecolor{pear}{rgb}{0.82, 0.89, 0.19}
\definecolor{pyellow}{rgb}{1.0, 0.97, 0.12}
\definecolor{lgreen}{rgb}{0.75, 1.0, 0.0}
\definecolor{brown}{rgb}{0.6, 0.4, 0.08}
\definecolor{ggray}{rgb}{0.27, 0.35, 0.27}
\definecolor{darkgreen}{HTML}{00AA00}

%%%%%%%%%%%%%%%%%%%%%%%%%%%%%%%%%%%%%%%%

\DeclareRobustCommand{\SkipTocEntry}[5]{}


\ctikzset{bipoles/resistor/height=0.15}
\ctikzset{bipoles/resistor/width=0.4}
\usetikzlibrary {positioning}
\usetikzlibrary{patterns}
\pgfplotsset{compat=1.10}
\usepgfplotslibrary{fillbetween}

\usetikzlibrary{patterns}
\usetikzlibrary{calc,decorations.markings}
\usetikzlibrary{shapes,snakes}


\tikzstyle{Cwhite}=[scale =.8,circle, fill = white, minimum size=3mm]
\tikzstyle{Cgray}=[scale =.4,circle, fill = gray, minimum size=3mm]
\tikzstyle{Cblack2}=[scale =.4,circle, fill = black, minimum size=5mm]
\tikzstyle{Cblack}=[scale =.7,circle, fill = black, minimum size=3mm]
\tikzstyle{C0}=[scale =.9,circle, fill = black!0, inner sep = 0pt, minimum size=3mm]
\tikzstyle{C1}=[scale =.7,circle, fill = black!0, inner sep = 0pt, minimum size=3mm]
\tikzstyle{Cred}=[scale =.4,circle, fill = red, minimum size=3mm]


\newcommand{\symbuptm}{\tikz[scale=.37, baseline=-0.7]{\draw(.1,-.05)to[out=50,in=-130](.9,.05) (.5,0)--++(0,.5);}}
%\newcommand{\symbuptm}{\tikz[scale=.28, baseline=-0.7]{\draw(.1,-.05)to[out=50,in=-130](.9,.05) (.5,0)--++(0,.5);}}
%\newcommand{\almorth}{\tikz[scale=.28, baseline=-0.7]{\symbuptm \draw (-.44,.22) circle (14pt);}}
\newcommand{\almorth}{\tikz[scale=.37, baseline=-0.7]{\symbuptm \draw (-.44,.22) circle (14pt);}}
\newcommand{\aplus}{\hspace{-.2cm}\almorth\,}
\newcommand{\aperp}{\,\symbuptm\,}

%%%%%%%%%%%%%%%%%%%%%%%%%%%%%%%%%%%%%%%%%%%%%%%%%%%%%%%

%%%%%%%%%%%%%%%%%%%%%%%%%%%%%%%%%%%%%%%%%%%%%%%%%%%%%%%%%%\langleetc modifiés
\newcommand{\innone}[2]{{\langle#2\rangle_{#1}}}

\newcommand{\binnone}[2]{{\left\langle#2\right\rangle_{#1}}}

\newcommand{\llangle}{\mathopen{\langle\mkern -3.1mu\lvert}}
\newcommand{\rrangle}{\mathclose{\rvert\mkern -3.1mu\rangle}}

\newcommand{\bllangle}{\mathopen{\bigl\langle\mkern -3.1mu\bigl\lvert}}
\newcommand{\brrangle}{\mathclose{\bigr\rvert\mkern -3.1mu\bigr\rangle}}

\newcommand{\highinn}[2]{{\llangle#2\rrangle_{#1}}}
\newcommand{\bhighinn}[2]{{\bllangle#2\brrangle_{#1}}}
\newcommand{\highinnoneorth}[2]{{\llangle#2\rrangle^{\perp}_{#1}}}
\newcommand{\innoneorth}[2]{{\langle#2\rangle^{\perp}_{#1}}}



\newcommand{\symbupl}{\tikz[scale=.1, baseline=-1.7]{\draw(.98,-.75)to[out=90,in=90](.98,.86) ;}}
\newcommand{\langleup}{\langle \tikz[scale=.28, baseline=-0.7]{\symbupl}}
\newcommand{\symbupr}{\tikz[scale=.1, baseline=-1.7]{\draw(.91,-.77)to[out=90,in=90](.91,.86) ;}}
\newcommand{\rangleup}{\rangle\tikz[scale=.28, baseline=-0.7]{\symbupr}}

\newcommand{\highinnprime}[2]{{\llangle #2\rrangle'_{#1}}}
%%%%%%%%%%%%%%%%%%%%%%%%%%%%%%%%%%%%%%%%%%%%%%%MODIF DE \bar pour overbar
\makeatletter
\newsavebox\myboxA
\newsavebox\myboxB
\newlength\mylenA

\newcommand*\overbar[2][0.75]{
\sbox{\myboxA}{$\m@th#2$}
\setbox\myboxB\null
\ht\myboxB=\ht\myboxA
\dp\myboxB=\dp\myboxA
\wd\myboxB=#1\wd\myboxA
\sbox\myboxB{$\m@th\overline{\copy\myboxB}$}
\setlength\mylenA{\the\wd\myboxA}
\addtolength\mylenA{-\the\wd\myboxB}
\ifdim\wd\myboxB<\wd\myboxA
\rlap{\hskip 1\mylenA\usebox\myboxB}{\usebox\myboxA}
\else
\hskip -0.5\mylenA\rlap{\usebox\myboxA}{\hskip 0.5\mylenA\usebox\myboxB}
\fi}
\makeatother
\newcommand{\grind}[1]{#1}
\newcommand{\comp}[1]{\overbar[.5]{#1}}
%%%%%
\newcommand{\Hyb}{{\mathit{hyb}}}
\newcommand{\hyb}{^{\Hyb}}

\newcommand{\mg}{\mathscr M}
\newcommand{\mgbar}{\comp{\mathscr M}}
\newcommand{\mgg}[1]{\mg_{#1}^{ }}
\newcommand{\mgbarg}[1]{\mgbar_{#1}^{ }}
\newcommand{\mghyb}[1]{\mg_{#1}^{\hyb}}

%%%%%%%%%%%%%%%%%%%%%%%%%%%%%%%%%%%%%%%%%%%%%%%%%%%%%%%
\newcommand{\norm}[1]{\|#1\|}
\newcommand{\normsq}[1]{\|#1\|^{2}}
\newcommand{\normbsq}[1]{\left\|#1\right\|^{2}}
\newcommand{\refnorm}[1]{\lvert#1\rvert}
\newcommand{\refnormsq}[1]{\lvert#1\rvert^{2}}
\newcommand{\refnormbsq}[1]{\left\lvert#1\right\rvert^{2}}
\newcommand{\normorth}[1]{\|#1\|^{\perp}}

\newcommand{\rest}[1]{\vert_{#1}}

\newcommand{\abs}[1]{\lvert #1\rvert}
\newcommand{\abss}[1]{\left\lvert #1\right\rvert}



\newcommand{\rquot}[2]{#1/\,#2}
\newcommand{\st}{\bigm|}

\newcommand{\orient}{\mathcal{O}}
%%%%%%%%%%%%%%%%%%%%%%%%%%%%%%%%%%%%%%%%%%%%
\DeclareMathOperator{\rk}{rank}
\DeclareMathOperator{\Hom}{Hom}
\DeclareMathOperator{\id}{Id}
\DeclareMathOperator{\weak}{w}
\DeclareMathOperator{\projeuc}{p}
\DeclareMathOperator{\Jac}{Jac}
\DeclareMathOperator{\vol}{vol}
\DeclareMathOperator{\Vor}{Vor}

\DeclareMathOperator{\spann}{span}

\renewcommand{\Re}{\operatorname{Re}}
\DeclareMathOperator{\intt}{int}
\DeclareMathOperator{\dist}{d}
\newcommand{\graphgenus}{{h}}
\newcommand{\filter}{\mathrm{F}}
\newcommand{\Q}{\mathbb{Q}}
\newcommand{\R}{\mathbb{R}}
\newcommand{\Z}{\mathbb{Z}}
\newcommand{\N}{\mathbb{N}}
\newcommand{\C}{\mathbb{C}}
\newcommand{\E}{\mathbb{E}}
\newcommand{\T}{\mathbb{T}}
\renewcommand{\L}{\mathbb{L}}
\renewcommand{\S}{\mathbb{S}}

\newcommand{\smallcc}{\mathbb{C}}
\newcommand{\mgr}{\mathcal{G}}
\newcommand{\curve}{\mathcal{C}}
\newcommand{\rsf}{\mathcal{S}}
\newcommand{\Ical}{\mathcal{I}}

\newcommand{\grm}[2]{\mathrm{gr}_{#1}^{#2}}


\newcommand{\proj}{\mathfrak{p}}
\newcommand{\qf}{\mathfrak{q}}
\newcommand{\fin}{\mathfrak f}
\newcommand{\scS}{\mathscr{S}}

\newcommand{\transpose}{\mathrm{T}}

\renewcommand{\setminus}{\smallsetminus}
\renewcommand{\emptyset}{\varnothing}
\newcommand{\dtor}{\zeta}

\newcommand{\interior}[1]{\intt(#1)}






\newcommand{\hdist}{d_{\mathrm{H}}}
\newcommand{\GH}{\mathrm{GH}}
\newcommand{\HA}{\mathrm{H}}



\newcommand{\suppp}[1]{{\mathrm{supp}(#1)^+}}
\newcommand{\suppm}[1]{{\mathrm{supp}(#1)^-}}
\newcommand{\supppm}[1]{{\mathrm{supp}(#1)^\pm}}



\newcommand{\lexeq}{\preceq_{\mathrm{lex}}}
\newcommand{\lexst}{\prec_{\mathrm{lex}}}
\newcommand{\gexeq}{\succeq_{\mathrm{lex}}}
\newcommand{\gex}{\succ_{\mathrm{lex}}}



%%%%%%%%%%%%%%%%%%%%%%%%%%%%%%%%%%%%%%%%%%%%%%%%%
\renewcommand*{\to}{\mathchoice{\longrightarrow}{\rightarrow}{\rightarrow}{\rightarrow}}
%%%%%%%%%%
\renewcommand*{\implies}{\mathchoice{\Longrightarrow}{\Rightarrow}{\Rightarrow}{\Rightarrow}}
%%%%%%%%%%
\let\oldmapsto\mapsto
\renewcommand*{\mapsto}{\mathchoice{\longmapsto}{\oldmapsto}{\oldmapsto}{\oldmapsto}}
%%%%%%%%%%%
\newcommand{\longhookrightarrow}{\lhook\joinrel\longrightarrow}
\newcommand*{\hookto}{\mathchoice{\longhookrightarrow}{\hookrightarrow}{\hookrightarrow}{\hookrightarrow}}
%%%%%
%Better widebar

\makeatletter
\let\save@mathaccent\mathaccent
\newcommand*\if@single[3]{%
  \setbox0\hbox{${\mathaccent"0362{#1}}^H$}%
  \setbox2\hbox{${\mathaccent"0362{\kern0pt#1}}^H$}%
  \ifdim\ht0=\ht2 #3\else #2\fi
  }
%The bar will be moved to the right by a half of \macc@kerna, which is computed by amsmath:
\newcommand*\rel@kern[1]{\kern#1\dimexpr\macc@kerna}
%If there's a superscript following the bar, then no negative kern may follow the bar;
%an additional {} makes sure that the superscript is high enough in this case:
\newcommand*\widebar[1]{\@ifnextchar^{{\wide@bar{#1}{0}}}{\wide@bar{#1}{1}}}
%Use a separate algorithm for single symbols:
\newcommand*\wide@bar[2]{\if@single{#1}{\wide@bar@{#1}{#2}{1}}{\wide@bar@{#1}{#2}{2}}}
\newcommand*\wide@bar@[3]{%
  \begingroup
  \def\mathaccent##1##2{%
%Enable nesting of accents:
    \let\mathaccent\save@mathaccent
%If there's more than a single symbol, use the first character instead (see below):
    \if#32 \let\macc@nucleus\first@char \fi
%Determine the italic correction:
    \setbox\z@\hbox{$\macc@style{\macc@nucleus}_{}$}%
    \setbox\tw@\hbox{$\macc@style{\macc@nucleus}{}_{}$}%
    \dimen@\wd\tw@
    \advance\dimen@-\wd\z@
%Now \dimen@ is the italic correction of the symbol.
    \divide\dimen@ 3
    \@tempdima\wd\tw@
    \advance\@tempdima-\scriptspace
%Now \@tempdima is the width of the symbol.
    \divide\@tempdima 10
    \advance\dimen@-\@tempdima
%Now \dimen@ = (italic correction / 3) - (Breite / 10)
    \ifdim\dimen@>\z@ \dimen@0pt\fi
%The bar will be shortened in the case \dimen@<0 !
    \rel@kern{0.6}\kern-\dimen@
    \if#31
      \overline{\rel@kern{-0.6}\kern\dimen@\macc@nucleus\rel@kern{0.4}\kern\dimen@}%
      \advance\dimen@0.4\dimexpr\macc@kerna
%Place the combined final kern (-\dimen@) if it is >0 or if a superscript follows:
      \let\final@kern#2%
      \ifdim\dimen@<\z@ \let\final@kern1\fi
      \if\final@kern1 \kern-\dimen@\fi
    \else
      \overline{\rel@kern{-0.6}\kern\dimen@#1}%
    \fi
  }%
  \macc@depth\@ne
  \let\math@bgroup\@empty \let\math@egroup\macc@set@skewchar
  \mathsurround\z@ \frozen@everymath{\mathgroup\macc@group\relax}%
  \macc@set@skewchar\relax
  \let\mathaccentV\macc@nested@a
%The following initialises \macc@kerna and calls \mathaccent:
  \if#31
    \macc@nested@a\relax111{#1}%
  \else
%If the argument consists of more than one symbol, and if the first token is
%a letter, use that letter for the computations:
    \def\gobble@till@marker##1\endmarker{}%
    \futurelet\first@char\gobble@till@marker#1\endmarker
    \ifcat\noexpand\first@char A\else
      \def\first@char{}%
    \fi
    \macc@nested@a\relax111{\first@char}%
  \fi
  \endgroup
}
\makeatother

\let\oldbar\bar
\renewcommand*{\bar}[1]{{\mathchoice{\widebar{#1}}{\widebar{#1}}{\widebar{#1}}{\oldbar{#1}}}}

\newcommand*{\mk}{\mkern -1mu}
\newcommand*{\Mk}{\mkern -2mu}
\newcommand*{\mK}{\mkern 1mu}
\newcommand*{\MK}{\mkern 2mu}

\newcommand*{\relabel}{\renewcommand{\labelenumi}{(\theenumi)}}
\newcommand*{\romanenumi}{\renewcommand*{\theenumi}{\roman{enumi}}\relabel}
\newcommand*{\Romanenumi}{\renewcommand*{\theenumi}{\Roman{enumi}}\relabel}
\newcommand*{\alphenumi}{\renewcommand*{\theenumi}{\alph{enumi}}\relabel}
\newcommand*{\alphebisnumi}{\renewcommand*{\theenumi}{\alph{enumi}$'$}\relabel}
\newcommand*{\Alphenumi}{\renewcommand*{\theenumi}{\Alph{enumi}}\relabel}
\let\oldtilde\tilde
\renewcommand*{\tilde}[1]{\mathchoice{\widetilde{#1}}{\widetilde{#1}}{\oldtilde{#1}}{\oldtilde{#1}}}
\let\oldhat\hat
\renewcommand*{\hat}[1]{\mathchoice{\widehat{#1}}{\widehat{#1}}{\oldhat{#1}}{\oldhat{#1}}}
\let\oldexists\exists
\renewcommand*{\exists}{\mathrel{\oldexists}}
\let\oldforall\forall
\renewcommand*{\forall}{\mathrel{\oldforall}}

\newtheorem{cdrthm}{Theorem}[section]
\newtheorem{theo}[cdrthm]{Theorem}
\newtheorem{lemm}[cdrthm]{Lemma}
\newtheorem{prop}[cdrthm]{Proposition}
\theoremstyle{definition}\newtheorem{defi}[cdrthm]{Definition}
\newtheorem{rema}[cdrthm]{Remark}
\newtheorem{exam}[cdrthm]{Example}
\newtheorem{coro}[cdrthm]{Corollary}
\title[Almost orthogonal decomposition of Voronoi closures]{Higher rank inner products, Voronoi tilings and metric degenerations of tori: original Theorem4.1}
\author{Omid Amini and Noema Nicolussi}
\date{Source: Annales Henri Lebesgue8(2025),1109--1188; DOI10.5802/ahl.257}
\begin{document}\maketitle
\noindent\textit{Editorial scope.} The counted unit is original Theorem4.1 with its complete proof and the required source-local core: higher rank inner products, the induced filtration and graded pieces, almost orthogonality and the Lifting Lemma, the Voronoi cell definition, the basic properties of higher rank Voronoi cells, admissible discrete sets and the Voronoi decomposition theorem. Illustrated examples, the pullback/tameness machinery and every later application are omitted. The two figure references and the section reference in Remark4.2 are kept as external links with their published numbers, since the artwork itself is not redistributed; all mathematics is native editable text and nothing was written or repaired.
\section{Original basic notations}
\subsection*{Basic notations} \label{ss:BasicNotations}
For a point $a$ in the vector space $\R^r$, we denote by $a_1, \dots, a_r$ the coordinates of~$a$. We set $\R_+ \coloneqq (0, +\infty)$ and denote by $\R_+^r =(0,+\infty)^r$ the set of $r$-vectors with positive coordinates.

Let $H$ be a finite dimensional real vector space. In this paper, a polyhedron is a subset $P \subset H$ which can be written as an intersection of finitely many half-spaces in $H$. A polytope is a bounded polyhedron. A generalized polyhedron is a subset $P \subset H$ such that $P \cap Q$ is a polytope for every polytope $Q \subset H$. The terminology is borrowed from~\cite[p.~250]{Gruber}.


A (generalized) polyhedral tiling of a vector space $H$ in this paper means a (possibly infinite) collection $\Sigma$ of full dimensional (generalized) polyhedra $\sigma \subseteq H$ which verifies the following properties:
\begin{itemize}
\item The union of $\sigma \in \Sigma$ is $H$.
\item The interiors of $\sigma \in \Sigma$ are disjoint open sets in $H$.
\end{itemize}

Note that this definition allows two top-dimensional tiles to intersect only on part of their proper faces, see Figure~\href{https://ahl.centre-mersenne.org/articles/10.5802/ahl.257/}{4.1} for an example.


We use the notation $\highinn{}{\cdot\,,\cdot}$ for inner products with values in $\R^r$, $r\in \N$. For scalar products, we usually use $\innone{}{\cdot\,,\cdot}$.

\section{Higher rank inner products} \label{sec:HigherRankInnerProducts}

\subsection{Totally ordered vector spaces}



Let $(\Lambda, \preceq)$ be a totally ordered real vector space of finite dimension. It is known that $(\Lambda, \preceq)$ is isomorphic to $\R^r$ with its lexicographic order $\lexeq$ where $r$ is the dimension of $\Lambda$~\cite{Bir-Lattice, HW52}. This means there is no loss of generality in assuming that $\Lambda = \R^r$ and $\preceq =\lexeq$. For $a =(a_1, \dots, a_r)$ and $b=(b_1, \dots, b_r)$ in $\R^r$, we have $a \preceq b$ if either $a=b$, or there exists $j\in [r]$ such that $a_1=b_1, \dots, a_{j-1}=b_{j-1}$ and $a_j<b_j$.


In this paper, we introduce and study inner products with values in $(\R^r, \preceq=\lexeq)$. The notion can be generalized to handle inner products with values in general totally ordered vector spaces. However, upon fixing an isomorphism, we can always reduce to the above case, which has the advantage of simplifying the notation. For this reason, we treat inner products with values in $(\R^r, \preceq=\lexeq)$, and discuss the general case in Section~\href{https://ahl.centre-mersenne.org/articles/10.5802/ahl.257/}{11.6}.



The vector space $\Lambda = \R^r$ with its lexicographic order $\preceq=\lexeq$ has a natural decreasing filtration
\begin{equation}\label{eq:FiltrationRr}
\Lambda^\bullet \colon \, \Lambda^1 \coloneq \Lambda\supset \Lambda^2 \supset \dots \supset \Lambda^{r} \supset \Lambda^{r+1} \coloneq (0)
\end{equation}
defined by
\[
\Lambda^j \coloneqq \bigl\{a = (a_1, \dots, a_r) \in \R^r \bigm| a_1 = \dots =a_{j-1}=0\bigr\}.
\]
Note that $\Lambda^j$ is naturally isomorphic to $\R^{r+1-j}$ endowed with its lexicographic order.




\subsection{Inner product spaces of higher rank}

Let $H$ be a real vector space of finite dimension and $\Lambda =\R^r$ with its lexicographical order $\preceq=\lexeq$.

Consider a symmetric bilinear form $\highinn{}{\cdot \,, \cdot} \colon H \times H \to \R^r$. We denote by $\highinn{j}{\cdot\,, \cdot}$ the $j^{\rm th}$ coordinate of $\highinn{}{\cdot\,,\cdot}$, so that $\highinn{}{\cdot\,,\cdot}$ has the form
\[
\highinn{}{\cdot\,,\cdot} = \bigl(\highinn{1}{\cdot\,, \cdot}, \dots, \highinn{r}{\cdot\,, \cdot}\bigr).
\]
In case that $r=1$, we call $\highinn{}{\cdot \,, \cdot}$ a \emph{scalar form}.

We say that $\highinn{}{\cdot \,, \cdot}$ is \emph{non-negative} if $\highinn{}{x\,, x} \succeq 0$ for all $x \in H$, and \emph{positive} if $\highinn{}{x\,, x} \succ 0$ for all $x \in H \setminus \{0\}$. In the case $r=1$, a positive scalar form is called a \emph{scalar product}.

A symmetric bilinear form $\highinn{}{\cdot \,, \cdot} \colon H \times H \to \R^r$ defines a filtration on $H$ as follows. Consider the decreasing filtration $\Lambda^\bullet$ on $\Lambda = \R^r$ given in~\eqref{eq:FiltrationRr}. Setting
\begin{align}
\filter^j \,\MK &\!\Mk:= \bigl\{ x \in H \bigm| \,\highinn{}{x \,, y} \in \Lambda^j \text{ for all }y \in H\bigr\}\\
&= \bigl\{ x \in H \bigm| \,\highinn{}{x \,, y}_i = 0 \quad\text{for all }i <j\text{  and all }y \in H\bigr\}
\end{align}
we obtain a non-increasing filtration
\[
\filter^\bullet\colon \filter^1 = H \supseteq \filter^2 \supseteq \filter^3 \supseteq \dots \supseteq \filter^r \supset \filter^{r+1} = (0)
\]
of $H$. We call $\filter^\bullet$ the \emph{filtration induced by} $\highinn{}{\cdot \,, \cdot}$.

\setcounter{cdrthm}{3}
The following properties are immediate from the definition of the filtration $\filter^\bullet$.

\begin{prop}\label{prop:GradedForms}
Let $\highinn{}{\cdot \,, \cdot} \colon H \times H \to \R^r$ be a symmetric bilinear form. For $j \in [r]$, the $j^{\rm th}$ component $\highinn{j}{\cdot \,, \cdot}$ of $\highinn{}{\cdot \,, \cdot}$ induces a symmetric bilinear form
\[
\highinn{j}{\cdot\,, \cdot} \colon \rquot{\filter^j}{\filter^{j+1}} \times \rquot{\filter^j}{\filter^{j+1}} \to \R
\]
on the quotient $\rquot{\filter^j}{\filter^{j+1}}$.

If $\highinn{}{\cdot \,, \cdot}$ is non-negative on $H$, then $\highinn{j}{\cdot \,, \cdot}$ is non-negative on $\rquot{\filter^j}{\filter^{j+1}}$ for all $j$. If $\highinn{j}{\cdot \,, \cdot}$ is positive on $\rquot{\filter^j}{\filter^{j+1}}$ for all $j$, then $\highinn{}{\cdot \,, \cdot}$ is positive on $H$.
\end{prop}

\setcounter{cdrthm}{4}
\begin{defi}[Inner product of higher rank] \label{defi:InnerProduct}\leavevmode
\begin{itemize}
\item An \emph{inner product on $H$ with values in $\Lambda = \R^r$} is a symmetric bilinear form
\[
\highinn{}{\cdot\,, \cdot} \colon H \times H \to \R^r
\]
such that for all $j \in [r]$, the induced form $\highinn{j}{\cdot\,, \cdot}$ on $\filter^j / \filter^{j+1}$ is positive, equivalently, if $\highinn{j}{\gamma, \gamma} >0$ for all $\gamma \in \filter^j \setminus \filter^{j+1}$. The pair $(H, \highinn{}{\cdot\,, \cdot})$ is called \emph{an inner product space}.

\item For $j \in [r]$, the $j^{\rm th}$ \emph{graded piece} $\grm{}{j}H$ of $(H, \highinn{}{\cdot\,, \cdot})$ is the quotient
\[
\grm{}{j}H = \rquot{\filter^j}{\filter^{j+1}}
\]
endowed with the scalar product
\[
\highinn{j}{\cdot\,, \cdot} \colon \grm{}{j}H\times \grm{}{j}H \to \R.
\]
We denote by $\proj_j \colon \filter^j \to \grm{}{j}H$ the projection map.

\item By a slight abuse of notation, we refer to inner products with values in $\R^r$, for arbitrary positive integer $r$, as \emph{inner products of higher rank} or simply as \emph{inner products} (although, strictly speaking, $r=1$ is allowed and, for larger~$r$, their image is allowed to be a one-dimensional subspace of $\R^r$).%\qedd
\end{itemize}
\end{defi}

\setcounter{cdrthm}{16}
\subsection{The almost orthogonal decomposition} \label{ss:AlmostOrthogonalDecomposition}
Let $\highinn{}{\cdot\,,\cdot} \colon H \times H \to \R^r$ be an inner product and $\filter^\bullet$ the induced filtration on $H$.



Consider two elements $\gamma \in \filter^{j} \setminus \filter^{j+1}$ and $\gamma' \in \filter^{j'} \setminus \filter^{j'+1}$. Then $\highinn{}{\gamma\,,\gamma'}$ lies in ${\Lambda}^{\max\{j, j'\}}$, that is, it has the form
\[
\highinn{}{\gamma\,,\gamma'} = \bigl(0, \dots, 0, \highinn{\max\{j, j'\}}{\gamma\,,\gamma'}, \dots, \highinn{r}{\gamma\,,\gamma'} \bigr).
\]
This motivates the following definition.

\begin{defi}[Almost orthogonality]
Two elements $\gamma, \gamma' \in H$ are called \emph{almost orthogonal} and we write $\gamma \aperp\gamma'$ if $\highinn{\max\{j,j'\}}{\gamma, \gamma'} = 0$ where $j, j' \in [r]$ are the indices with $\gamma \in \filter^j \setminus \filter^{j+1}$ and $\gamma' \in \filter^{j'} \setminus \filter^{j'+1}$.

Two subsets $A, B \subset E$ are called \emph{almost orthogonal} and we write $A \aperp B$, if any pair of elements $\gamma \in A$ and $\eta \in B$ are almost orthogonal. We use the notation $A \aplus B$ for $A+B$, indicating that $A$ and $B$ are almost orthogonal.%\qedd
\end{defi}

The terminology is justified as follows.

\begin{lemm}\label{lem:orthogonality_vanishing}
If $\gamma \in H$ is almost orthogonal to itself, then $\gamma = 0$. In particular, if $U, V \subset H$ are almost orthogonal subspaces, then $U \aplus V$ is a direct sum. That is, if $u+v = 0$ for $u\in U$ and $v \in V$, then $u=v = 0$.
\end{lemm}

\begin{proof}
If $\gamma \neq 0$, then there exists $k \in [r]$ with $\gamma \in \filter^k \smallsetminus \filter^{k+1}$. Moreover, $\highinn{k}{\gamma, \gamma} >0$ since $\highinn{}{\cdot, \cdot}$ is an inner product. On the other hand, if $\gamma$ is orthogonal to itself, then $\highinn{k}{\gamma, \gamma} =0$. We have arrived at a contradiction.
\end{proof}


As we discuss next, the notion of almost orthogonality allows to canonically decompose $H$ in terms of its graded pieces $\grm{}{j}H$, $j \in [r]$.

\begin{lemm}[Lifting lemma] \label{lem:LiftingLemma}
Let $\highinn{}{\cdot\,,\cdot} \colon H \times H \to \R^r$ be an inner product. Then for $j\in [r]$, there exists a unique linear map $\proj_j^* \colon \grm{}{j}H \to \filter^j$ with the following properties:
\begin{enumerate}
\item $\proj_j \circ \proj_j^* = \id$ and, in particular, $\proj_j^* \colon \grm{}{j}H \hookto \filter^j$ is an embedding.
\item The image of $\proj_j^*$ is almost orthogonal to $\filter^{j+1}$.
\end{enumerate}
The inner product space $(H, \highinn{}{\cdot, \cdot})$ has the almost orthogonal decomposition
\[
H = \proj_1^*\left(\grm{}{1}H\right) \aplus\,\proj_2^*\left(\grm{}{2}H\right) \aplus \cdots\!\aplus\, \proj_r^*\left(\grm{}{r}H\right).
\]
That is, $H$ is the direct sum of the subspaces $\proj_j^*(\grm{}{j}(H))$, $j \in [r]$, and they are pairwise almost orthogonal.
\end{lemm}

\begin{defi}
We refer to the embedding $\proj_j^*\colon \grm{}{j} H \hookto H$ as the \emph{canonical lifting} of $\grm{}{j}H$ to $H$, and denote its image by $H_j \coloneqq \proj_j^*(\grm{}{j} H)$. The resulting almost orthogonal decomposition is denoted by
\[
H = H_1\aplus H_2 \aplus \cdots \aplus H_r. %\qedd
\]
\end{defi}

Note that the $j^{\rm th}$ component $\highinn{j}{\cdot \,, \cdot}$ of the inner product $\highinn{}{\cdot \,, \cdot}$ is a scalar product on $H_j$ for all $j\in [r]$. Namely, by the properties in Lemma~\ref{lem:LiftingLemma}, it is clear that $H_j \subset F^j$ and $H_j \aperp F^{j+1}$. Applying Lemma~\ref{lem:orthogonality_vanishing}, it follows that every $x \in H \setminus \{0\}$ belongs to $F^j \setminus F^{j+1}$ and consequently satisfies $\highinn{j}{x \,, x}>0$.


\begin{proof}[Proof of Lemma~\ref{lem:LiftingLemma}]
Fix $j \in [r]$ and let $\gamma\in \grm{}{j}H$. We show that there exists a unique element $\theta \in \filter^j$ such that $\proj_j(\theta) =\gamma$, and $\theta$ is almost orthogonal to $\filter^{j+1}$. This defines a unique linear map $\proj_j^*$ with the above properties.

To prove the uniqueness of $\theta$, let $\theta$ and $\theta'$ be two elements of $\filter^j$ such that $\proj_j(\theta) =\proj_j(\theta') = \gamma$, and both are almost orthogonal to $\filter^{j+1}$. Then the difference $\theta-\theta'$ belongs to $\filter^{j+1}$ and is moreover almost orthogonal to $\filter^{j+1}$. Lemma~\ref{lem:orthogonality_vanishing} then ensures that $\theta =\theta'$.


In order to prove the existence of $\theta$, we proceed inductively, and show that for each $k=0, \dots, r-j$, we can find an element $\theta_{j+k}\in \filter^j$ such that
\begin{enumerate}\romanenumi
\item\label{prflemm2.19.1} $\proj_j(\theta_{j+k}) = \gamma$, and

\item\label{prflemm2.19.2} for each $h=1, \dots, k$, we have $\highinn{j+h}{\theta_{j+k}, \eta} = 0 \text{ for all }\eta\in \filter^{j+h}.$
\end{enumerate}
Setting $\theta = \theta_{r}$, we then obtain the desired element.

For $k=0$, choose any $\theta_j \in \filter^j$ with $\proj_j(\theta_j) =\gamma$. Then $\theta_j$ verifies~\eqref{prflemm2.19.1} and~\eqref{prflemm2.19.2} is empty.

Assume we have found $\theta_{j+k}$ for some $k< r-j$. Define the linear form
\begin{align*}
f_{j+k}\colon \filter^{j+k+1} &\to \R\\
\eta &\longmapsto \highinn{j+k+1}{\theta_{j+k}, \eta}.
\end{align*}
Since $\highinn{j+k+1}{\cdot \,, \cdot}$ is an inner product on $\grm{}{j+k+1}H$, there exists an element $\alpha \in \filter^{j+k+1}$ with
\[
f_{j+k}(\eta) = \highinn{j+k+1}{\alpha, \eta}, \quad \forall \eta\in \filter^{j+k+1}.
\]
Let $\theta_{j+k+1} \coloneqq \theta_{j+k} - \alpha$ and note that $\proj_j(\theta_{j+k+1}) = \proj_j(\theta_{j+k}) =\gamma$. Since $\alpha \in \filter^{j+k+1}$, we also have $\highinn{j+h}{\alpha, \eta}=0$ for any $\eta \in \filter^{j+h}$ for all $h=1, \dots, k$. Finally,
\[
\highinn{j+k+1}{\theta_{j+k+1}, \eta} = \highinn{j+k+1}{\theta_{j+k}, \eta} - \highinn{j+k+1}{\alpha, \eta}= 0, \quad \forall \eta\in \filter^{j+k+1},
\]
by construction. Thus, $\theta_{j+k+1}$ satisfies the properties~\eqref{prflemm2.19.1} and~\eqref{prflemm2.19.2}. This finishes the proof of existence and uniqueness of the embedding $\proj_j^* \colon \grm{}{j}H \hookto H$.

The properties of $\proj_j^*$ clearly imply that the subspaces $\proj_i^*(\grm{}{j}H)$, $j \in [r]$, are pairwise almost orthogonal. It remains to prove that $H$ is the direct sum of these subspaces. Since $\dim(H) = \sum_{j} \dim(\grm{}{j}H)$, the desired claim follows from Lemma~\ref{lem:orthogonality_vanishing}.
\end{proof}

\section{Higher rank Voronoi decompositions} \label{sec:HigherRankVoronoiDecompositions}
Throughout this section, let $H$ be a finite dimensional vector space and 
let further $\highinn{}{\cdot\, ,\cdot}\colon H\times H \to \R^r$ be an inner product taking values in $\Lambda = \R^r$ with the lexicographic order $\preceq=\lexeq$.

Let $S \subset H$ be a discrete subset of $H$, that is, $S$ has no accumulation point in $H$ endowed with its Euclidean topology. We define the \emph{Voronoi cell} of $\gamma \in S$ as
\begin{equation}\label{eq:DefinitionVoronoiCell}
\Vor_S(\gamma) \coloneqq \bigl\{x\in H \bigm| \highinn{}{x-\gamma, x- \gamma} \preceq \highinn{}{x-\eta, x-\eta} \quad\text{for all } \eta \in S\bigr\}.
\end{equation}
If $S$ is understood from the context, we simply write $\Vor(\gamma)$.


The analogy with the case $r = 1$ and the relation to degenerating scalar products (see Theorem~\href{https://ahl.centre-mersenne.org/articles/10.5802/ahl.257/}{2.13}) suggest to interpret $\highinn{}{x, x}$ as the (vector-valued) length of an~element $x \in H$. The Voronoi cell of $\gamma$ is precisely the set of all points $x \in H$ whose closest point in $S$ is $\gamma$.


\subsection{Basic properties and pathologies} \label{ss:Pathologies}
In the classical setting of a scalar product, that is, $r=1$, the Voronoi cells $\Vor_S(\gamma)$, $\gamma \in S$, are convex, closed and have non-empty, mutually disjoint interiors. If $S \subset H$ is finite or $S= \L$ is a lattice in $H$, then each Voronoi cell is a polyhedron, that is, an intersection of finitely many half-spaces in $H$. For a general discrete subset $S \subset H$, the Voronoi cells are generalized polyhedra (see Section~\ref{ss:BasicNotations} for the terminology). Moreover, they provide a decomposition $H = \bigcup_{\gamma \in S} \Vor(\gamma)$ of $H$ into (generalized) polyhedra, which is called the \emph{Voronoi decomposition}, see e.g.~\cite[p.~465]{Gruber}.

\setcounter{cdrthm}{3}
The following proposition summarizes the basic properties of higher rank Voronoi cells.

\begin{prop}\label{prop:general_voronoi}
Let $\highinn{}{\cdot\,,\cdot}\colon H\times H \to \R^r$ be an inner product and $S$ a discrete subset of $H$. Then for any $\gamma\in S$, the Voronoi cell $\Vor(\gamma)$ is a non-empty convex set. Moreover, the interiors of the Voronoi cells are mutually disjoint.

In general, $\Vor(\gamma)$ is not closed and can have empty interior. The union of the Voronoi cells can be a proper subset of $H$. Moreover, even in case that $S = \L$ is a lattice of full rank in $H$, the Voronoi cell $\Vor(\gamma)$ can be unbounded.
\end{prop}



Suppose that $0 \in S$. A point $x$ belongs to the Voronoi cell $\Vor(0)$ of $0 \in S$ exactly when
\[
\highinn{}{x,x} \preceq \highinn{}{x - \eta, x-\eta} \quad \text{for all $\eta \in S$}.
\]
For $\eta \in S$, the above inequality can be written as
\[
2 \highinn{}{x,\eta} \preceq \highinn{}{\eta,\eta}.
\]
This formulation involves the $\R^r$-valued linear form $2 \highinn{}{\cdot,\eta}$ on $H$ and the constant $\highinn{}{\eta,\eta} \in \R^r$. The set of points $x \in H$ satisfying this inequality may be viewed as a higher rank version of a half-space in $H$.

\setcounter{cdrthm}{4}
\subsection{Admissible discrete sets and their Voronoi decompositions} \label{ss:AdmissibleDiscreteSetsAndVoronoiDecomposition}
The above examples suggest that discrete sets in the Euclidean topology on $H$ are not necessarily well compatible with higher rank inner products. In this section, we formulate a notion of discreteness in the higher rank setting, and prove that this leads to well-behaving Voronoi decompositions.



\subsubsection{Admissible discrete sets} Let $\highinn{}{\cdot \,, \cdot} \colon H \times H \to \R^r$ be an inner product with values in $\R^r$. Let further $H = H_1 \aplus \cdots \aplus H_r$ be the almost orthogonal decomposition. We write elements $x \in H$ as $x=x_1+\dots+x_r$ with $x_j \in H_j$, $j= 1, \dots, r$.

Consider a subset $S\subset H$. For each $j \in [r]$ and a tuple $z = (z_1, \dots, z_{j-1})$, with $z_i \in H_i$ for $i=1,\dots, j-1$, define the subset $S_{/z} \subseteq S$ as
\begin{equation}\label{eq:SOver}
S_{/z} \coloneqq \left\{ \gamma=\gamma_1 + \dots + \gamma_r \in S \st \gamma_i = z_i \text{ for $i=1, \dots, j-1$} \right\}.
\end{equation}
Moreover, we define the translated set
\begin{equation}\label{eq:Translation}
TS(z) \coloneqq S_{/z} - z_1 - \dots -z_{j-1}.
\end{equation}
We have $TS(z) \subset \filter^j$. In the special case $j=1$, we have the empty vector $z = \varnothing$ and recover
\[
S_{/\varnothing} = TS (\varnothing) = S \subset \filter^1.
\]
\begin{defi}[Admissible discrete sets]
A subset $S \subset H$ is called \emph{admissible discrete} for the inner product $\highinn{}{\cdot\,,\cdot}$, if for all $j \in [r]$ and all tuples $z= (z_1, \dots, z_{j-1})$, with $z_i \in H_i$ for $i=1,\dots, j-1$, the projection $\proj_j(TS(z))$ is a discrete set in $\grm{}{j} H$.%\qedd
\end{defi}

\setcounter{cdrthm}{8}
\subsubsection{Voronoi decomposition induced by an admissible discrete set}

We have the following Voronoi decomposition theorem for admissible discrete sets.
\begin{theo}\label{thm:voronoi-discrete}
Let $\highinn{}{\cdot \,, \cdot} \colon H \times H \to \R^r$ be an inner product and $S \subset H$ an~admissible discrete set. Then, the sets $\Vor_S(\gamma)$, $\gamma\in S$, cover the full space, that is, $H = \bigcup_{\gamma \in S} \Vor_S(\gamma)$.
\end{theo}

We call the above decomposition $H = \bigcup_{\gamma \in S} \Vor_S(\gamma)$ the \emph{Voronoi decomposition} associated to the inner product $\highinn{}{\cdot \,, \cdot} $ and the admissible discrete set $S$.



In the proof, we use the almost orthogonal decomposition $H= \aplus_{j=1}^r H_j$ and decompose each element $x \in H$ as $x = x_1 + \dots + x_r$ with $x_j \in H_j$.

\section{Structure of Voronoi cells} \label{sec:StructureVoronoiCells}
Let $\highinn{}{\cdot \,, \cdot} \colon H \times H \to \R^r$ be an inner product taking values in $\Lambda = \R^r$ with lexicographic order $\preceq=\lexeq$ and $S \subset H$ an admissible discrete subset. Consider the associated Voronoi decomposition $\bigcup_{\gamma\in S} \Vor_S(\gamma)=H$ (see Theorem~\ref{thm:voronoi-discrete}). In this section, we provide a description of the closed Voronoi cells $\bar{\Vor}_{S}(\gamma)$.


Consider the filtration $\filter^\bullet = (\filter^j)_{j}$ on $H$ induced by $\highinn{}{\cdot\,,\cdot}$. Recall that the $j^{\rm th}$ graded piece $\grm{}{j}H = \rquot{\filter^j}{\filter^{j+1}}$, $j \in [r]$, carries the scalar product $\highinn{j}{\cdot\,,\cdot} \colon \grm{}{j}H \times \grm{}{j}H \to \R$. We have the almost orthogonal decomposition $H = \aplus_{j=1}^r H_j$. In the following, we decompose each $x \in H$ as $x = x_1 + \dots +x_r$ with $x_j \in H_j$, $j=1, \dots, r$. Note that $\highinn{j}{\cdot\,,\cdot}$ is a scalar product on $H_j$. Moreover, we have $H_j \subseteq \filter^j$ and $H_j \simeq \grm{}{j}H$, where the isomorphism is given by $\proj_j^*$ in the Lifting Lemma~\ref{lem:LiftingLemma}.

Consider an element $\gamma =\gamma_1 + \dots + \gamma_r$ in $S$. For $j \in [r]$, define the subset $S_{\gamma, j} \subset H_j$ by
\[
S_{\gamma, j} = \left\{ \theta_j \bigm| \theta \in S \quad\text{and}\quad \theta_i = \gamma_i \;\text{for}\; i \in\{1, \dots, j-1\} \right\}.
\]
Note that $S_{\gamma, 1} = \{\theta_1 \mid \theta \in S\}$ does not depend on $\gamma$.

Using the tuple $(\gamma_1, \dots, \gamma_{j-1})$ and the notation from~\eqref{eq:SOver} and~\eqref{eq:Translation}, we have
\[
S_{\gamma, j} = \left\{ \theta_j \st \theta \in S_{/(\gamma_1, \dots, \gamma_{j-1})} \right\} = \proj_j^\ast\bigl(\proj_j(TS(\gamma_1, \dots, \gamma_{j-1}))\bigr).
\]
Since $S$ is admissible discrete, $S_{\gamma,j}$ is a discrete subset of $H_j$ with $\gamma_j \in S_{\gamma,j}$. Let
\[
V_{\gamma, j} \coloneqq \Vor_{S_{\gamma,j}}(\gamma_j) \subseteq H_j
\]
be the Voronoi cell of $\gamma_j$ in $H_j$ associated to $S_{\gamma,j}$ and the scalar product $\highinn{j}{\cdot\,,\cdot}$.



The main result of this section reads as follows.


\begin{theo}[Almost orthogonal decomposition theorem for Voronoi cells]\label{thm:VoronoiCellDecomposition-discrete}
Notations as above, let $S \subset H$ be an admissible discrete set and $\gamma \in S$. Then, the closure of $\Vor_{S}(\gamma)$ has the almost orthogonal decomposition
\[
\bar\Vor_{S}(\gamma) = V_{\gamma,1} \aplus \cdots \aplus V_{\gamma,r}.
\]
That is, $\bar\Vor_{S}(\gamma)$ coincides with the Minkowski sum of the $V_{\gamma,j}$, $j \in [r]$, and these sets are pairwise almost orthogonal.

It follows that $\bar\Vor_{S}(\gamma)$ is a (generalized) polyhedron with a non-empty interior.
\end{theo}

\begin{rema}
We stress that taking the closure of the Voronoi cell in Theorem~\ref{thm:VoronoiCellDecomposition-discrete} is necessary, since the Voronoi cells are not always closed (see Figure~\href{https://ahl.centre-mersenne.org/articles/10.5802/ahl.257/}{4.1}, as well as Figure~\href{https://ahl.centre-mersenne.org/articles/10.5802/ahl.257/}{10.2} and Section~\href{https://ahl.centre-mersenne.org/articles/10.5802/ahl.257/}{10} for examples).%\qedd
\end{rema}

\begin{proof}[Proof of Theorem~\ref{thm:VoronoiCellDecomposition-discrete}] We first prove the claim in the case $0\in S$ and $\gamma=0$, and then explain how this implies the general case. In this case, we have
\[
S_j \coloneqq S_{\gamma,j}= \left\{\theta_j \bigm| \theta \in S\quad\text{and}\quad\theta_1=\dots = \theta_{j-1}=0 \right\} = \left\{\theta_j \bigm| \theta \in S \cap \filter^j\right\}
\]
and
\[
V_j \coloneqq V_{\gamma,j} = \Vor_{S_j}(0), \quad j=1,\dots, r.
\]
Since $V_{j} \subset H_j$, we have $V_i \aperp V_j$ for $i\neq j$. We will show that $\bar \Vor_{S}(0) = V_1 \aplus \cdots \aplus V_r$.

We first prove the inclusion $\bar \Vor_{S}(0) \subseteq V_1 \aplus \dots \aplus V_r$. It will suffice to show that
\[
\Vor_{S}(0) \subseteq V_1 \aplus \cdots \aplus V_r,
\]
as the right hand side is clearly closed in $H$. Let $x\in \Vor_{S}(0)$. Then we have the inequalities
\[
2 \highinn{}{x, \gamma} \preceq \highinn{}{\gamma, \gamma} \quad \forall\, \gamma \in S.
\]
Looking at the first coordinate in both sides, we obtain the inequalities
\[
2 \highinn{1}{x, \gamma} \leq \highinn{1}{\gamma,\gamma}\,\quad \forall \gamma \in S.
\]
Using the almost orthogonal decomposition, these inequalities can be rewritten as
\[
2\highinn{1}{x_1, \gamma_1} \leq \highinn{1}{\gamma_1,\gamma_1}\,\quad \forall \gamma \in S.
\]
This shows that $x_1$ belongs to the Voronoi cell $V_1=\Vor_{S_1}(0)$ in $H_1$.

The difference $x-x_1$ belongs to $\filter^2$, and for all $\gamma \in \filter^2 \cap S$, we have
\[
2\highinn{2}{x-x_1,\gamma} = 2\highinn{2}{x,\gamma} \preceq \highinn{2}{\gamma,\gamma},
\]
which is equivalent to the inequalities
\[
2\highinn{2}{x_2,\gamma_2} \preceq \highinn{2}{\gamma_2,\gamma_2} \quad \forall \gamma \in \filter^2 \cap S.
\]
This means that $x_2$ belongs to the Voronoi cell $V_2 = \Vor_{S_2}(0)$ in $H_2$. Proceeding inductively in $j$, we find that $x_1 \in V_1, x_2\in V_2, \dots, x_j\in V_j$, and
\[
2\bhighinn{j+1}{x_{j+1},\gamma_{j+1}} \preceq \bhighinn{j+1}{\gamma_{j+1},\gamma_{j+1}} \quad \forall \gamma \in \filter^{j+1} \cap S.
\]
For $j=r$, we obtain $x \in V_1\aplus \dots \aplus V_r$, as required.

We now show the reverse inclusion $V_1 \aplus \cdots \aplus V_r \subseteq \bar \Vor_{S}(0)$. Since the interior $\interior{V_j}$ is dense in $V_j$ (viewed as a subset of $H_j$), it will be enough to prove the inclusion
\[
\mathring V_1 \aplus \cdots \aplus \mathring V_r \subseteq \Vor_{S}(0)
\]
for the sets
\[
\mathring V_j \coloneqq \interior{V_j} = \intt\left(\Vor_{S_j}(0)\right).
\]
More precisely, here we take the interior of $V_j$ as a subset of $H_j$, and view it as a subset of $H$. Let $x =x_1+ \dots +x_r$ be a point in the left hand side, with $x_j \in \mathring V_j$. We prove that $x \in \Vor_{S}(0)$ by verifying the inequalities
\[
2\highinn{}{x, \gamma} \preceq \highinn{}{\gamma, \gamma} \quad \forall \gamma \in S.
\]
We will in fact prove the strict inequalities.

Take an element $\gamma \in S \setminus \{0\}$, and let $j\in [r]$ be the integer with $\gamma \in \filter^j \smallsetminus \filter^{j+1}$. Then,
\[
\highinn{1}{\gamma,\gamma}=\dots = \highinn{j-1}{\gamma,\gamma}=0, \quad \highinn{j}{\gamma, \gamma}>0.
\]
The Lifting Lemma~\ref{lem:LiftingLemma} implies that
\[
\highinn{j}{x_i, \gamma} =0 \quad \text{ for } i \neq j.
\]
For $i= j$, we have
\[
\highinn{j}{x_j, \gamma} = \highinn{j}{x_j, \gamma_j}.
\]
Since $x_j \in \mathring V_j$, and $\gamma \in S\cap \filter^{j} \smallsetminus \filter^{j+1}$, it follows that $\gamma_j \neq 0$ and
\[
2\highinn{j}{x_j, \gamma_j} < \highinn{j}{\gamma_j, \gamma_j} = \highinn{j}{\gamma, \gamma}.
\]
Combining the above, we infer that, as required,
\[
2\highinn{}{x, \gamma} = \left(0, \dots, 0, 2\highinn{j}{x_j, \gamma}, \ast, \dots, \ast\right) \prec \left(0, \dots, 0, \highinn{j}{\gamma, \gamma}, \ast, \dots, \ast\right) = \highinn{}{\gamma, \gamma}.
\]



The remaining claims on $\bar \Vor_{S}(0)$ follow from the decomposition just established and the properties of the (generalized) polyhedra $V_j$, $j=1, \dots, r$.


The statement for $\gamma \in S$ follows by applying the decomposition to the translated set $S' = S-\gamma$, which is again admissible discrete, and noticing that
\begin{align*}
\Vor_S(\gamma) &= \Vor_{S'}(0)+\gamma\\
 &=\left(\Vor_{S'_1}(0) +\gamma_1\right) \aplus \cdots \aplus \left(\Vor_{S'_r}(0)+\gamma_r\right) = V_{\gamma, 1}\aplus \cdots \aplus V_{\gamma, r}. \qedhere
\end{align*}
\end{proof}

\begin{thebibliography}{99}
\bibitem[Bir40]{Bir-Lattice} Garrett Birkhoff, \emph{Lattice theory}, Colloquium Publications25, American Mathematical Society,1940.
\bibitem[HW52]{HW52} Melvin Hausner and James G.Wendel, ``Ordered vector spaces'', \emph{Proc. Am. Math. Soc.}\textbf{3}(1952),977--982.
\bibitem[Gru07]{Gruber} Peter M.Gruber, \emph{Convex and Discrete Geometry}, Grundlehren der Mathematischen Wissenschaften336, Springer,2007.
\end{thebibliography}
\end{document}
